\begin{center}
Lecture Notes

on

Elliptic Differential Equations

Qing Han

Fang-Hua Lin
\end{center}

\begin{center}
Contents
\end{center}

Preface \hfill v

Chapter 1.\ \ Harmonic Functions \hfill 1
\[
\begin{array}{llr}
1.1.&\text{Mean Value Properties}&1\\
1.2.&\text{Fundamental Solutions}&7\\
1.3.&\text{Maximum Principles}&14\\
1.4.&\text{Energy Methods}&19
\end{array}
\]

Chapter 2.\ \ Maximum Principles \hfill 23
\[
\begin{array}{llr}
2.1.&\text{Strong Maximum Principle}&23\\
2.2.&\text{Poisson Equations}&28\\
2.3.&\text{A Priori Estimates}&36\\
2.4.&\text{Gradient Estimates}&38\\
2.5.&\text{Alexandroff Maximum Principle}&42\\
2.6.&\text{The Moving Plane Method}&48
\end{array}
\]

Chapter 3.\ \ Weak Solutions, Part I \hfill 51
\[
\begin{array}{llr}
3.1.&\text{Growth of Local Integrals}&52\\
3.2.&\text{H\"older Continuity of Solutions}&58\\
3.3.&\text{H\"older Continuity of Gradients}&62
\end{array}
\]

Chapter 4.\ \ Weak Solutions, Part II \hfill 67
\[
\begin{array}{llr}
4.1.&\text{Local Boundedness}&67\\
4.2.&\text{The H\"older Continuity}&78\\
4.3.&\text{Moser's Harnack Inequality}&82\\
4.4.&\text{Nonlinear Equations}&91
\end{array}
\]

Chapter 5.\ \ Viscosity Solutions \hfill 95
\[
\begin{array}{llr}
5.1.&\text{Alexandroff Maximum Principle}&95\\
5.2.&\text{The Harnack Inequality}&100\\
5.3.&\text{Schauder Estimates}&108\\
5.4.&W^{2,p}\text{ Estimates}&112
\end{array}
\]

Bibliography \hfill 117

Index \hfill 119

\chapter*{Preface}

The purpose of these note is to present some basic methods for obtaining various
a priori estimates for second order partial differential equations of elliptic type
with particular emphasis on maximum principles, Harnack inequalities, and their
applications. The equations we deal with are always linear, although most of the
methods obviously apply to nonlinear problems. Students with some knowledge of
real variables and Sobolev functions should be able to follow these notes without
much difficulty.

It is not our intention to give a complete account of the related theory. Our
goal is simply to provide these notes as a bridge between the books of John [4] and
of Evans [2], which also study equations of other types, and the book of Gilbarg
and Trudinger [3] which gives a relatively complete account of the theory of elliptic
equations of second order. We also hope our notes can serve as a bridge between
the book of Krylov [5] on the classical theory of elliptic equations developed before
or around the 1960s and the book by Caffarelli and Cabré [1] which studies fully
nonlinear elliptic equations, the theory obtained in the 1980s.

\hfill Qing Han, Fang-Hua Lin
